\usepackage{graphicx}
\usepackage{hyperref}

\usepackage[utf8]{inputenc}
\usepackage[T1]{fontenc}
\usepackage{amsmath,amssymb} 
\usepackage{enumitem}
\begin{document}

\hypersetup{
    pdftitle={Laboratorio 3 ICN-292},
    pdfauthor={Isidora Cordova Stevenson},
    colorlinks=true,
    urlcolor=blue
}

\begin{titlepage}
\centering
\includegraphics[scale=.3]{Imagenes/Logo}
\vspace{12mm}

\textsc{\Large Sistema de Información para la Gestión }\\[0.5cm]
\textit{\large
\\2do. Semestre 2026\hspace{5mm}\textbar\hspace{5mm}
Programa Diurno \hspace{5mm}\textbar\hspace{5mm}
\ Campus Vitacura}\\[0.5cm]
\vspace{8mm}
\HRule \\[6mm]
{\huge \ \textbf{Informe de Laboratorio 4 } }\\[0.4cm]
\HRule \\[10mm]

\begin{minipage}[t]{0.5\textwidth}
\begin{flushleft} \large
\emph{\textbf{Autora:}}\\
         Isidora Renata  \textsc{Córdova Stevenson} \\
	{\small 202204631-7, Ingeniería Civil Industrial}\\
\end{flushleft}
\end{minipage}
~ % Important!
\begin{minipage}[t]{0.4\textwidth}
\begin{flushright} \large
\emph{\textbf{Profesores:}} \\
Jose Miguel  \textsc{González P.
} \\
Jose Luis  \textsc{Saéz T.
} \\
\end{flushright}
\end{minipage}\\[12mm]

{\large \textbf{\textit{Enlace Github:}}}\\[1.5mm]
{\small \url{https://github.com/Isiusm/ICN292-Lab4-Cordova-Isidora}}\\[5mm]

\vfill
{\large \ Campus Vitacura}\\[5mm]
\end{titlepage}

\setlength{\parskip}{2mm}

\setlength{\parskip}{4mm}
\setlength{\parindent}{0cm}
\newpage



\clearpage
\section*{Modelo entidad-relación}

\begin{center}
    \includegraphics[
        width=\textwidth,
        height=0.65\textheight,
        keepaspectratio
    ]{\detokenize{ICN292-Lab4-Córdova-Isidora-EER.png}}
\end{center}



\clearpage
\section*{Consulta 1}

\textbf{Código SQL}

\begin{lstlisting}
SELECT nombre, genero
FROM autores
WHERE pais <> 'Chile' OR pais IS NULL
ORDER BY nombre ASC;
\end{lstlisting}

La comparación \texttt{pais <> 'Chile'} devuelve los países distintos de Chile, pero no incluye los valores \texttt{NULL}, porque SQL no considera verdadera una comparación con un valor nulo. Por eso se agrega \texttt{OR pais IS NULL}.

\evidencia{Consulta 1.png}{Resultado de la consulta 1.}



\clearpage
\section*{Consulta 2}

\textbf{Código SQL}

\begin{lstlisting}
SELECT titulo, paginas
FROM libros
WHERE paginas BETWEEN 190 AND 220
ORDER BY paginas ASC;
\end{lstlisting}

\evidencia{Consulta 2.png}{Resultado de la consulta 2.}


\clearpage
\section*{Consulta 3}

\textbf{Código SQL}

\begin{lstlisting}
SELECT
    COUNT(*) AS total_libros,
    SUM(ejemplares_vendidos) AS total_vendidos,
    AVG(paginas) AS promedio_paginas,
    MIN(anio) AS anio_mas_antiguo,
    MAX(ejemplares_vendidos) AS mayor_venta
FROM libros;
\end{lstlisting}

\evidencia{Consulta 3.png}{Resultado de la consulta 3.}


\clearpage
\section*{Consulta 4}

\textbf{Código SQL}

\begin{lstlisting}
SELECT DISTINCT genero
FROM autores
ORDER BY genero ASC;
\end{lstlisting}

\evidencia{Consulta 4.png}{Resultado de la consulta 4.}


\clearpage
\section*{Consulta 5}

\textbf{Código SQL}

\begin{lstlisting}
SELECT
    a.nombre,
    COUNT(*) AS cantidad_libros
FROM autores a
JOIN libros l
    ON a.id_autor = l.id_autor
WHERE l.anio = 2024
GROUP BY a.id_autor, a.nombre
HAVING COUNT(*) > 1
ORDER BY a.nombre ASC;
\end{lstlisting}

\evidencia{Consulta 5.png}{Resultado de la consulta 5.}


\clearpage
\section*{Consulta 6}

\textbf{Código SQL}

\begin{lstlisting}
SELECT titulo, ejemplares_vendidos
FROM libros
WHERE anio = 2024
AND ejemplares_vendidos > 2000
ORDER BY ejemplares_vendidos DESC
LIMIT 3;
\end{lstlisting}

\evidencia{Consulta 6.png}{Resultado de la consulta 6.}


\clearpage
\section*{Consulta 7}

\textbf{Código SQL}

\begin{lstlisting}
INSERT INTO libros
(id_autor, id_coleccion, titulo, paginas,
 ejemplares_vendidos, anio)
VALUES
(3, NULL, 'Papel Mojado', 260, 50, 2024);

SELECT *
FROM libros
WHERE titulo = 'Papel Mojado';

DELETE FROM libros
WHERE titulo = 'Papel Mojado';

SELECT *
FROM libros
WHERE titulo = 'Papel Mojado';
\end{lstlisting}

\textbf{Evidencias de ejecución}

\evidencia[0.32\textheight]
{Captura 7 (1).png}
{Registro del libro antes de eliminarlo.}

\evidencia[0.32\textheight]
{Captura 7 (2eliminado).png}
{Comprobación posterior al DELETE.}


\clearpage
\section*{Consulta 8}

\textbf{Código SQL}

\begin{lstlisting}
SELECT titulo, ejemplares_vendidos
FROM libros
WHERE titulo = 'Cielo de Enero';

UPDATE libros
SET ejemplares_vendidos = ejemplares_vendidos + 100
WHERE titulo = 'Cielo de Enero';

SELECT titulo, ejemplares_vendidos
FROM libros
WHERE titulo = 'Cielo de Enero';
\end{lstlisting}

\textbf{Evidencias de ejecución}

\evidencia[0.34\textheight]
{Consulta 8 Update (1).png}
{Primera evidencia de la consulta 8.}

\evidencia[0.34\textheight]
{Consulta 8 Update (2).png}
{Segunda evidencia de la consulta 8.}


\clearpage
\section*{Consulta 9}

\textbf{Código SQL}

\begin{lstlisting}
INSERT INTO colecciones (titulo, fecha_publicacion)
VALUES
('Voces del Sur', '2020-03-15'),
('Horizontes Literarios', '2022-08-20');

UPDATE libros
SET id_coleccion = 1
WHERE id_libro BETWEEN 1 AND 6;

UPDATE libros
SET id_coleccion = 2
WHERE id_libro BETWEEN 7 AND 12;

SELECT
    libros.titulo AS titulo_libro,
    colecciones.titulo AS titulo_coleccion,
    colecciones.fecha_publicacion
FROM libros
INNER JOIN colecciones
    ON libros.id_coleccion = colecciones.id_coleccion
ORDER BY libros.id_libro;
\end{lstlisting}

\evidencia{Consulta 9.png}{Resultado de la consulta 9.}

\clearpage
\section*{Consulta 10}

\textbf{Código SQL}

\begin{lstlisting}
INSERT INTO pagos (id_autor, fecha_pago, monto)
VALUES
(1, '2026-01-15', 150000.00),
(1, '2026-06-15', 200000.00),
(2, '2026-03-10', 175000.00);

SELECT
    autores.nombre,
    pagos.fecha_pago,
    pagos.monto
FROM pagos
INNER JOIN autores
    ON pagos.id_autor = autores.id_autor
ORDER BY autores.nombre;

SELECT
    autores.nombre,
    SUM(pagos.monto) AS total_pagado
FROM autores
INNER JOIN pagos
    ON autores.id_autor = pagos.id_autor
GROUP BY autores.id_autor, autores.nombre
ORDER BY autores.nombre;
\end{lstlisting}

\textbf{Evidencias de ejecución}


\evidencia[0.34\textheight]
{Captura 10 .png}
{Detalle de los pagos realizados a los autores.}

\evidencia[0.34\textheight]
{Captura 10 (2).png}
{Total pagado a cada autor.}
\end{document}


